\documentclass[11pt]{article}
\usepackage{../../../tex/lecturestyle}

\title{Prime Sieves}
\author{}
\date{}

\begin{document}
\thispagestyle{fancy}

\lecturetitleband{Number Theory}{Prime Sieves}{3}

\section{Overview}
Testing one integer for primality is cheap. Testing \emph{every} integer up to
$N$ one at a time is not: trial division alone costs $\Theta(\sqrt{n})$ per
number. The insight behind sieving is to invert the question. Instead of
asking ``is $n$ prime?'' for each $n$, we start from the primes we have
already found and strike out their multiples wholesale. Each composite is
then eliminated by the primes dividing it, and whatever survives is prime.

This lecture covers:
\begin{itemize}[leftmargin=1.5em]
  \item the \textbf{sieve of Eratosthenes}, enumerating all primes $\le N$ in
        $O(N \log\log N)$ time;
  \item the \textbf{smallest prime factor table}, which turns factorization of
        any $n \le N$ into an $O(\log n)$ lookup loop;
  \item the \textbf{linear sieve}, which touches each composite exactly once
        and so runs in $O(N)$;
  \item bulk evaluation of \textbf{multiplicative functions}
        ($d(n)$, $\sigma(n)$, $\varphi(n)$) over $1 \ldots N$;
  \item the \textbf{segmented sieve}, for primes in a narrow window $[L, R]$
        where $R$ is far too large to sieve from $1$.
\end{itemize}

\section{The Sieve of Eratosthenes}

\subsection{The algorithm}
Keep a boolean array $\mathrm{isPrime}[0 \ldots N]$, initially all true except
indices $0$ and $1$. Walk $n$ upwards. Whenever $\mathrm{isPrime}[n]$ is still
true, $n$ has survived every earlier strike-out, so it is prime --- and we
then mark $2n, 3n, 4n, \ldots$ as composite. By the time the walk reaches $n$,
every composite below $n$ has already been struck out by one of its prime
divisors, so the test ``still true'' is exactly the test ``prime''.

\begin{theorembox}[(Correctness)]
After the sweep, $\mathrm{isPrime}[n]$ is true if and only if $n$ is prime.

\emph{Proof sketch.} If $n$ is prime, nothing can mark it: a mark on index $n$
is only ever written while processing some $m$ with $2 \le m < n$ and $m \mid
n$, and a prime has no such divisor. Conversely if $n \ge 2$ is composite,
write $n = mp$ with $p$ its smallest prime factor; then $p < n$, and
$\mathrm{isPrime}[p]$ was still true when the walk reached $p$, so the inner
loop for $p$ ran and marked the multiple $n$. \qed
\end{theorembox}

% FIGURE: sieve_grid
\begin{figure}[H]
\centering
\begin{tikzpicture}[x=0.9cm,y=0.78cm]
\node[lcell,minimum width=0.85cm,minimum height=0.72cm,lgray,text=lectureink!40] at (0,0) {1};
\node[lcell,minimum width=0.85cm,minimum height=0.72cm,fill=white,very thick,draw=lectureink] at (1,0) {\textbf{2}};
\node[lcell,minimum width=0.85cm,minimum height=0.72cm,fill=white,very thick,draw=lectureink] at (2,0) {\textbf{3}};
\node[lcell,minimum width=0.85cm,minimum height=0.72cm,fill=lectureaccent!30,text=lectureink!70] at (3,0) {4};
\node[lcell,minimum width=0.85cm,minimum height=0.72cm,fill=white,very thick,draw=lectureink] at (4,0) {\textbf{5}};
\node[lcell,minimum width=0.85cm,minimum height=0.72cm,fill=lectureaccent!30,text=lectureink!70] at (5,0) {6};
\node[lcell,minimum width=0.85cm,minimum height=0.72cm,fill=white,very thick,draw=lectureink] at (6,0) {\textbf{7}};
\node[lcell,minimum width=0.85cm,minimum height=0.72cm,fill=lectureaccent!30,text=lectureink!70] at (7,0) {8};
\node[lcell,minimum width=0.85cm,minimum height=0.72cm,fill=lecturewarn!30,text=lectureink!70] at (8,0) {9};
\node[lcell,minimum width=0.85cm,minimum height=0.72cm,fill=lectureaccent!30,text=lectureink!70] at (9,0) {10};
\node[lcell,minimum width=0.85cm,minimum height=0.72cm,fill=white,very thick,draw=lectureink] at (0,-1) {\textbf{11}};
\node[lcell,minimum width=0.85cm,minimum height=0.72cm,fill=lectureaccent!30,text=lectureink!70] at (1,-1) {12};
\node[lcell,minimum width=0.85cm,minimum height=0.72cm,fill=white,very thick,draw=lectureink] at (2,-1) {\textbf{13}};
\node[lcell,minimum width=0.85cm,minimum height=0.72cm,fill=lectureaccent!30,text=lectureink!70] at (3,-1) {14};
\node[lcell,minimum width=0.85cm,minimum height=0.72cm,fill=lecturewarn!30,text=lectureink!70] at (4,-1) {15};
\node[lcell,minimum width=0.85cm,minimum height=0.72cm,fill=lectureaccent!30,text=lectureink!70] at (5,-1) {16};
\node[lcell,minimum width=0.85cm,minimum height=0.72cm,fill=white,very thick,draw=lectureink] at (6,-1) {\textbf{17}};
\node[lcell,minimum width=0.85cm,minimum height=0.72cm,fill=lectureaccent!30,text=lectureink!70] at (7,-1) {18};
\node[lcell,minimum width=0.85cm,minimum height=0.72cm,fill=white,very thick,draw=lectureink] at (8,-1) {\textbf{19}};
\node[lcell,minimum width=0.85cm,minimum height=0.72cm,fill=lectureaccent!30,text=lectureink!70] at (9,-1) {20};
\node[lcell,minimum width=0.85cm,minimum height=0.72cm,fill=lecturewarn!30,text=lectureink!70] at (0,-2) {21};
\node[lcell,minimum width=0.85cm,minimum height=0.72cm,fill=lectureaccent!30,text=lectureink!70] at (1,-2) {22};
\node[lcell,minimum width=0.85cm,minimum height=0.72cm,fill=white,very thick,draw=lectureink] at (2,-2) {\textbf{23}};
\node[lcell,minimum width=0.85cm,minimum height=0.72cm,fill=lectureaccent!30,text=lectureink!70] at (3,-2) {24};
\node[lcell,minimum width=0.85cm,minimum height=0.72cm,fill=green!35!lectureink!12,text=lectureink!70] at (4,-2) {25};
\node[lcell,minimum width=0.85cm,minimum height=0.72cm,fill=lectureaccent!30,text=lectureink!70] at (5,-2) {26};
\node[lcell,minimum width=0.85cm,minimum height=0.72cm,fill=lecturewarn!30,text=lectureink!70] at (6,-2) {27};
\node[lcell,minimum width=0.85cm,minimum height=0.72cm,fill=lectureaccent!30,text=lectureink!70] at (7,-2) {28};
\node[lcell,minimum width=0.85cm,minimum height=0.72cm,fill=white,very thick,draw=lectureink] at (8,-2) {\textbf{29}};
\node[lcell,minimum width=0.85cm,minimum height=0.72cm,fill=lectureaccent!30,text=lectureink!70] at (9,-2) {30};
\node[lcell,minimum width=0.85cm,minimum height=0.72cm,fill=white,very thick,draw=lectureink] at (0,-3) {\textbf{31}};
\node[lcell,minimum width=0.85cm,minimum height=0.72cm,fill=lectureaccent!30,text=lectureink!70] at (1,-3) {32};
\node[lcell,minimum width=0.85cm,minimum height=0.72cm,fill=lecturewarn!30,text=lectureink!70] at (2,-3) {33};
\node[lcell,minimum width=0.85cm,minimum height=0.72cm,fill=lectureaccent!30,text=lectureink!70] at (3,-3) {34};
\node[lcell,minimum width=0.85cm,minimum height=0.72cm,fill=green!35!lectureink!12,text=lectureink!70] at (4,-3) {35};
\node[lcell,minimum width=0.85cm,minimum height=0.72cm,fill=lectureaccent!30,text=lectureink!70] at (5,-3) {36};
\node[lcell,minimum width=0.85cm,minimum height=0.72cm,fill=white,very thick,draw=lectureink] at (6,-3) {\textbf{37}};
\node[lcell,minimum width=0.85cm,minimum height=0.72cm,fill=lectureaccent!30,text=lectureink!70] at (7,-3) {38};
\node[lcell,minimum width=0.85cm,minimum height=0.72cm,fill=lecturewarn!30,text=lectureink!70] at (8,-3) {39};
\node[lcell,minimum width=0.85cm,minimum height=0.72cm,fill=lectureaccent!30,text=lectureink!70] at (9,-3) {40};
\node[lcell,minimum width=0.85cm,minimum height=0.72cm,fill=white,very thick,draw=lectureink] at (0,-4) {\textbf{41}};
\node[lcell,minimum width=0.85cm,minimum height=0.72cm,fill=lectureaccent!30,text=lectureink!70] at (1,-4) {42};
\node[lcell,minimum width=0.85cm,minimum height=0.72cm,fill=white,very thick,draw=lectureink] at (2,-4) {\textbf{43}};
\node[lcell,minimum width=0.85cm,minimum height=0.72cm,fill=lectureaccent!30,text=lectureink!70] at (3,-4) {44};
\node[lcell,minimum width=0.85cm,minimum height=0.72cm,fill=lecturewarn!30,text=lectureink!70] at (4,-4) {45};
\node[lcell,minimum width=0.85cm,minimum height=0.72cm,fill=lectureaccent!30,text=lectureink!70] at (5,-4) {46};
\node[lcell,minimum width=0.85cm,minimum height=0.72cm,fill=white,very thick,draw=lectureink] at (6,-4) {\textbf{47}};
\node[lcell,minimum width=0.85cm,minimum height=0.72cm,fill=lectureaccent!30,text=lectureink!70] at (7,-4) {48};
\node[lcell,minimum width=0.85cm,minimum height=0.72cm,fill=violet!25,text=lectureink!70] at (8,-4) {49};
\node[lcell,minimum width=0.85cm,minimum height=0.72cm,fill=lectureaccent!30,text=lectureink!70] at (9,-4) {50};
\node[llabel,anchor=west] at (10.0,-0.2) {crossed out first by:};
\node[lcell,minimum width=0.6cm,minimum height=0.5cm,fill=lectureaccent!30] at (10.4,-0.90) {}; \node[llabel,anchor=west] at (10.75,-0.90) {multiples of 2};
\node[lcell,minimum width=0.6cm,minimum height=0.5cm,fill=lecturewarn!30] at (10.4,-1.55) {}; \node[llabel,anchor=west] at (10.75,-1.55) {multiples of 3};
\node[lcell,minimum width=0.6cm,minimum height=0.5cm,fill=green!35!lectureink!12] at (10.4,-2.20) {}; \node[llabel,anchor=west] at (10.75,-2.20) {multiples of 5};
\node[lcell,minimum width=0.6cm,minimum height=0.5cm,fill=violet!25] at (10.4,-2.85) {}; \node[llabel,anchor=west] at (10.75,-2.85) {multiples of 7};
\node[lcell,minimum width=0.6cm,minimum height=0.5cm,very thick,draw=lectureink,fill=white] at (10.4,-3.5) {}; \node[llabel,anchor=west] at (10.75,-3.5) {prime (never crossed out)};
\end{tikzpicture}
\caption{The sieve of Eratosthenes on $1,\dots,50$. Only the primes $2,3,5,7$ up to $\lfloor\sqrt{50}\rfloor$ need to be used as sieving primes; each composite is coloured by the first (smallest) prime that crosses it out, and the numbers left unmarked are the 15 primes up to 50.}\label{fig:sieve_grid}
\end{figure}

\subsection{Complexity}
The inner loop for a prime $p$ performs $\lfloor N/p \rfloor$ writes, so the
total work is
\[
  O\!\left(N + \sum_{p \le N} \frac{N}{p}\right)
  \;=\; O\!\left(N \sum_{p \le N} \frac{1}{p}\right).
\]
The sum of reciprocals of primes grows like $\log\log N$ --- this is Mertens'
second theorem, $\sum_{p \le N} 1/p = \log\log N + M + o(1)$ --- giving a total
of
\[
  O(N \log\log N).
\]
For practical purposes $\log\log N$ is a small constant (under $4$ for
$N = 10^9$), so the sieve behaves close to linearly.

\begin{notebox}
It is a common slip to quote the sieve as $O(N \log N)$. That would be the
bound if we summed $N/k$ over \emph{all} $k \le N$ (the harmonic series).
Restricting to primes is what buys the improvement from $\log N$ to
$\log\log N$.
\end{notebox}

\subsection{A worthwhile constant-factor refinement}
Two easy observations shrink the work substantially:

\begin{enumerate}[leftmargin=1.5em]
  \item \textbf{Start striking at $p^2$, not $2p$.} Any composite multiple
        $kp$ with $k < p$ has a prime factor smaller than $p$, so it was
        already struck out during an earlier round. The first multiple that
        can still be unmarked is $p \cdot p$.
  \item \textbf{Only sieve with $p \le \sqrt{N}$.} If $p > \sqrt{N}$ then
        $p^2 > N$, so by the previous point there is nothing left in range to
        strike. The outer loop may stop as soon as $p^2 > N$.
\end{enumerate}

Neither changes the asymptotic bound, but together they remove a large amount
of redundant writing, and the second removes the outer loop's tail entirely.

\section{Smallest Prime Factors}

\begin{definitionbox}
For an integer $n \ge 2$, the \emph{smallest prime factor} $\mathrm{spf}(n)$ is
the least prime dividing $n$. A prime $p$ has $\mathrm{spf}(p) = p$.
\end{definitionbox}

A sieve can record more than a single bit. If, instead of writing ``false'',
each strike records \emph{which} prime did the striking, the array becomes a
factorization oracle. Care is needed to keep the \emph{smallest} such prime:
either write only when the slot is still empty, or process primes in
decreasing order so that the smallest prime writes last.

Given the table, factorizing is a short loop --- repeatedly divide out
$\mathrm{spf}(n)$:
\[
  n \;\longrightarrow\; n / \mathrm{spf}(n) \;\longrightarrow\; \cdots
  \;\longrightarrow\; 1 .
\]
Each step strips one prime factor, so the loop runs $\Omega(n)$ times where
$\Omega(n)$ is the number of prime factors with multiplicity --- at most
$\log_2 n$, since every factor is $\ge 2$.

\begin{notebox}
Factorizing \emph{all} of $1, 2, \ldots, N$ this way costs
$O(N \log\log N)$ in total, not $O(N \log N)$. The total number of prime
factors with multiplicity is $\sum_p (\lfloor N/p \rfloor + \lfloor N/p^2
\rfloor + \cdots)$; the first terms contribute $N \log\log N + O(N)$ as above,
and the higher powers contribute only $N \sum_p \frac{1}{p(p-1)} = O(N)$.
\end{notebox}

\section{The Linear Sieve}

In the sieve of Eratosthenes a composite is visited once per distinct prime
factor: $30$ gets struck three times, by $2$, $3$ and $5$. If we arrange for
each composite to be struck \emph{exactly once} --- by its smallest prime
factor --- the total work drops to $O(N)$.

\subsection{The rule}
Walk $n$ from $2$ to $N$, maintaining the list of primes found so far and the
table $\mathrm{spf}$. At each $n$:
\begin{itemize}[leftmargin=1.5em]
  \item If $\mathrm{spf}(n)$ is still unset, nothing below $n$ divides $n$, so
        $n$ is prime: set $\mathrm{spf}(n) = n$ and append $n$ to the list.
  \item Then, for each known prime $p$ in increasing order with
        $p \le \mathrm{spf}(n)$ and $np \le N$, set $\mathrm{spf}(np) = p$.
\end{itemize}
The guard $p \le \mathrm{spf}(n)$ is the whole trick.

\begin{theorembox}[(Each composite is produced exactly once)]
Let $m \le N$ be composite, $p = \mathrm{spf}(m)$ and $n = m/p$. Then the pair
$(n, p)$ satisfies the guard, and it is the \emph{only} pair that generates
$m$.

\emph{Why the guard holds.} If $n$ had a prime factor $q < p$, then $q$ would
divide $m = np$ as well, contradicting $p = \mathrm{spf}(m)$. Hence
$p \le \mathrm{spf}(n)$ and the inner loop does reach $p$ while processing $n$.

\emph{Why it is unique.} Suppose the algorithm writes $m = n'p'$ for some
other pair. The guard forces $p' \le \mathrm{spf}(n')$, so $p'$ is no larger
than any prime factor of $n'$, making $p'$ the smallest prime factor of
$n'p' = m$. Thus $p' = p$, and then $n' = m/p' = n$. \qed
\end{theorembox}

Since each of the $O(N)$ composites receives exactly one write, and the prime
list scan is charged to those writes, the whole sieve runs in $O(N)$ time.

\begin{notebox}
Asymptotically superior does not mean faster in practice. At the sizes typical
in contests ($N \le 10^8$), the sieve of Eratosthenes often wins: it writes to
memory in long regular strides that the hardware prefetcher handles well,
while the linear sieve chases the prime list and writes to scattered indices,
and it must store an integer per slot rather than a bit. Prefer the linear
sieve when you genuinely need $\mathrm{spf}$ or a multiplicative function; a
plain bitset Eratosthenes is usually the faster way just to list primes.
\end{notebox}

\section{Multiplicative Functions in Bulk}

\begin{definitionbox}
$f$ defined on the positive integers is \emph{multiplicative} if $f(1) = 1$ and
$f(ab) = f(a)f(b)$ whenever $\gcd(a, b) = 1$. It is \emph{completely
multiplicative} if $f(ab) = f(a)f(b)$ for all $a, b$, coprime or not.
\end{definitionbox}

A multiplicative $f$ is determined by its values on prime powers. Splitting off
the full power of the smallest prime, $n = m \cdot p^{e}$ with
$p = \mathrm{spf}(n)$ and $p \nmid m$, the parts are coprime, so
\[
  f(n) \;=\; f(m) \cdot f(p^{e}),
\]
and since $m < n$ the value $f(m)$ is already known. One left-to-right pass
therefore fills the whole table, provided $f(p^e)$ is cheap to evaluate. The
total cost of peeling off the exponents is $O(N)$, because
$\sum_{n \le N} v_p(n) \le N + N\sum_p \frac{1}{p(p-1)} = O(N)$.

Three standard instantiations:

\begin{examplebox}[: divisor count $d(n)$]
The number of divisors is multiplicative with $d(p^e) = e + 1$ (choose the
exponent of $p$ freely between $0$ and $e$). So $d(n) = d(m)\,(e+1)$.
\end{examplebox}

\begin{examplebox}[: divisor sum $\sigma(n)$]
The sum of divisors is multiplicative with
$\sigma(p^e) = 1 + p + \cdots + p^{e}$, evaluable by Horner's rule in $O(e)$.
\end{examplebox}

\begin{examplebox}[: Euler's totient $\varphi(n)$]
$\varphi(n)$ counts the integers in $[1, n]$ coprime to $n$. It is
multiplicative, with $\varphi(p^e) = p^{e-1}(p-1)$: among $1 \ldots p^e$ exactly
the $p^{e-1}$ multiples of $p$ fail to be coprime.
\end{examplebox}

A completely multiplicative $f$ needs even less: $f(n) = f(n/p) \cdot f(p)$ for
\emph{any} prime factor $p$, with no exponent bookkeeping. For example
$f(n) = n^{k} \bmod M$ can be tabulated by computing $p^{k} \bmod M$ by fast
exponentiation at primes only, and a single multiplication at composites. With
$\pi(N) = O(N / \log N)$ primes below $N$, that is
$O\!\left(N + \frac{N \log k}{\log N}\right)$ overall, beating the naive
$O(N \log k)$.

\section{The Segmented Sieve}

Now suppose we want the primes in $[L, R]$ where $R$ may be as large as
$10^{12}$ but the window $R - L + 1$ is manageable (say $\le 10^7$). Sieving
from $1$ is hopeless --- we cannot allocate $10^{12}$ slots --- but we do not
need to.

The key fact is that any composite $x \le R$ has a prime factor at most
$\sqrt{x} \le \sqrt{R}$. So it suffices to:
\begin{enumerate}[leftmargin=1.5em]
  \item sieve the small primes up to $\lfloor \sqrt{R} \rfloor$ normally
        (only $10^6$ slots when $R = 10^{12}$);
  \item allocate one boolean per element of the window, indexing $x$ at
        position $x - L$;
  \item for each base prime $p$, strike its multiples inside the window,
        starting from the first multiple of $p$ that is $\ge L$ (and
        $\ge p^2$, since smaller multiples of $p$ are either out of range or
        already handled).
\end{enumerate}
Whatever survives in the window is prime. The cost is
$O\!\left(\sqrt{R}\log\log R + (R - L + 1)\log\log R\right)$, and the memory is
proportional to the window plus $\sqrt{R}$.

\begin{notebox}
The same scaffolding factorizes an entire window: instead of one bit per slot,
store the residual cofactor, dividing out each base prime as it strikes.
Whatever is left above $1$ at the end is a single large prime factor, since a
number $\le R$ has at most one prime factor exceeding $\sqrt{R}$.
\end{notebox}

\section{Complexity Summary}

\begin{center}
\begin{tabular}{lll}
\toprule
\textbf{Task} & \textbf{Time} & \textbf{Space} \\
\midrule
Primes up to $N$ (Eratosthenes)      & $O(N \log\log N)$ & $N$ bits \\
Primes and $\mathrm{spf}$ (linear)   & $O(N)$            & $N$ words \\
Factorize one $n$ given spf          & $O(\log n)$       & --- \\
Factorize all of $1 \ldots N$        & $O(N \log\log N)$ & --- \\
Multiplicative $f$ over $1 \ldots N$ & $O(N)$            & $N$ words \\
Primes in $[L, R]$ (segmented)       & $O((\sqrt R + R - L)\log\log R)$ & $\sqrt R + (R-L)$ \\
\bottomrule
\end{tabular}
\end{center}

\section{Reference Implementations}

\codenote
\subsection{Rust}
\corelisting{Rust}{example.rs}

\subsection{Python}
\corelisting{Python}{example.py}

\section{Worked Example}
\begin{examplebox}
Run the linear sieve to $N = 12$. Processing $n = 2$: $\mathrm{spf}(2)$ is
unset, so $2$ is prime; the inner loop sets $\mathrm{spf}(4) = 2$. At $n = 3$:
prime, and we set $\mathrm{spf}(6) = 2$ (from $p = 2 \le \mathrm{spf}(3) = 3$)
and $\mathrm{spf}(9) = 3$. At $n = 4$: $\mathrm{spf}(4) = 2$ is already set, so
$4$ is composite; the inner loop runs only for $p = 2$, since $p = 3$ exceeds
$\mathrm{spf}(4) = 2$ --- this is precisely the guard preventing $12$ from
being struck here as $4 \times 3$. Instead $12$ is struck later at $n = 6$,
$p = 2$, matching $\mathrm{spf}(12) = 2$.

Factorizing $360$ with the table: $\mathrm{spf}(360) = 2 \Rightarrow 180
\Rightarrow 90 \Rightarrow 45$; then $\mathrm{spf}(45) = 3 \Rightarrow 15
\Rightarrow 5$; then $5$ is prime. So $360 = 2^3 \cdot 3^2 \cdot 5$, and for
instance $d(360) = 4 \cdot 3 \cdot 2 = 24$ and
$\varphi(360) = 2^2(2{-}1) \cdot 3(3{-}1) \cdot (5{-}1) = 4 \cdot 6 \cdot 4 = 96$.
\end{examplebox}

\section{Further Reading}
\begin{itemize}[leftmargin=1.5em]
  \item AtCoder Algorithm Lectures index (source of this curriculum):
        \url{https://info.atcoder.jp/entry/algorithm_lectures/index}
  \item cp-algorithms, \emph{Sieve of Eratosthenes}:
        \url{https://cp-algorithms.com/algebra/sieve-of-eratosthenes.html}
  \item cp-algorithms, \emph{Linear Sieve}:
        \url{https://cp-algorithms.com/algebra/prime-sieve-linear.html}
  \item P. Pritchard, ``Linear Prime-Number Sieves: a Family Tree'',
        \emph{Science of Computer Programming} 9(1), 1987, pp. 17--35.
\end{itemize}

\end{document}
